\chapter{Maximum Principles}

\section*{1. Weak maximum principles for scalar equations}

\subsection*{1.1. The heat equation with a gradient term}

\begin{proposition}[Scalar maximum principle, first version]
\label{ch04-prop-scalar-max-principle-pointwise-bounds}
\source{chow-knopf-ricci-flow:page-0106}
\dcref{Proposition 4.1}
\group{weak-maximum-principles-scalar-equations}

Let \(u : \mathcal M^n \times [0,T] \to \mathbb R\) be a \(C^2\) solution to the heat equation
\[
\frac{\partial u}{\partial t} = \Delta_g u
\]
on a closed Riemannian manifold. If there are constants \(C_1 \le C_2\) such that
\[
C_1 \le u(x,0) \le C_2
\]
for all \(x \in \mathcal M^n\), then
\[
C_1 \le u(x,t) \le C_2
\]
for all \(x \in \mathcal M^n\) and \(t \in [0,T)\).
\end{proposition}

\begin{theorem}[Scalar maximum principle, second version]
\label{ch04-thm-scalar-max-principle-supersolution-lower-bound}
\source{chow-knopf-ricci-flow:page-0106, chow-knopf-ricci-flow:page-0107}
\dcref{Theorem 4.2}
\group{weak-maximum-principles-scalar-equations}

Let \(g(t) : t \in [0,T)\) be a \(1\)-parameter family of Riemannian metrics and \(X(t) : t \in [0,T)\) a \(1\)-parameter family of vector fields on a closed manifold \(\mathcal M^n\). Let \(u : \mathcal M^n \times [0,T) \to \mathbb R\) be a \(C^2\) function. Suppose that there exists \(\alpha \in \mathbb R\) such that \(u(x,0) \ge \alpha\) for all \(x \in \mathcal M^n\), and that \(u\) is a supersolution of the heat-type equation
\[
\frac{\partial v}{\partial t} = \Delta_{g(t)} v + \langle X, \nabla v \rangle
\]
at any \((x,t) \in \mathcal M^n \times [0,T)\) such that \(u(x,t) < \alpha\). Then \(u(x,t) \ge \alpha\) for all \(x \in \mathcal M^n\) and \(t \in [0,T)\).
\end{theorem}

\begin{proof}
\source{chow-knopf-ricci-flow:page-0107}
\sketch If \(H : \mathcal M^n \times [0,T) \to \mathbb R\) is a \(C^2\) function and \((x_0,t_0)\) is a point and time where \(H\) attains its minimum over \(\mathcal M^n \times [0,t_0]\), then
\[
\frac{\partial H}{\partial t}(x_0,t_0) \le 0,\qquad \nabla H(x_0,t_0)=0,\qquad \Delta H(x_0,t_0)\ge 0.
\]

Consider
\[
H(x,t) := [u(x,t)-\alpha]+\varepsilon t+\varepsilon
\]
for any \(\varepsilon>0\). Since \(H\ge \varepsilon>0\) at \(t=0\), it suffices to show that \(H>0\) on \(\mathcal M^n \times [0,T)\). Using the supersolution hypothesis, at any point where \(u<\alpha\) one has
\[
\frac{\partial H}{\partial t}\ge \Delta H + \langle X,\nabla H\rangle+\varepsilon.
\]
If \(H\le 0\) somewhere, then because \(\mathcal M^n\) is compact there is a first time \(t_0>0\) at which \(H\) vanishes at some \(x_0\). At \((x_0,t_0)\), the hypothesis \(H(x_0,t_0)=0\) implies \(u(x_0,t_0)<\alpha\), so the previous inequality applies. Combining it with the minimum-point inequalities gives
\[
0\ge \frac{\partial H}{\partial t}(x_0,t_0)\ge \Delta H(x_0,t_0)+\langle X,\nabla H\rangle(x_0,t_0)+\varepsilon\ge \varepsilon>0,
\]
a contradiction. Therefore \(H>0\), and letting \(\varepsilon \downarrow 0\) yields the claim.
\end{proof}

\subsection*{1.2. The heat equation with a linear reaction term}

\begin{proposition}[Scalar maximum principle, third version]
\label{ch04-prop-scalar-max-principle-linear-reaction}
\source{chow-knopf-ricci-flow:page-0108}
\dcref{Proposition 4.3}
\group{weak-maximum-principles-scalar-equations}

Let \(u : \mathcal M^n \times [0,T) \to \mathbb R\) be a \(C^2\) supersolution to
\[
\frac{\partial v}{\partial t} = \Delta_{g(t)} v + \langle X, \nabla v \rangle + \beta v
\]
on a closed manifold. Suppose that for each \(\tau \in [0,T)\) there exists \(C_\tau<\infty\) such that \(\beta(x,t)\le C_\tau\) for all \(x \in \mathcal M^n\) and \(t \in [0,\tau]\). If \(u(x,0)\ge 0\) for all \(x \in \mathcal M^n\), then \(u(x,t)\ge 0\) for all \(x \in \mathcal M^n\) and \(t \in [0,T)\).
\end{proposition}

\begin{proof}
\source{chow-knopf-ricci-flow:page-0108}
\sketch For a fixed \(\tau\in(0,T)\), set
\[
J(x,t)\doteq e^{-C_\tau t}u(x,t).
\]
Then
\[
\frac{\partial J}{\partial t}\ge \Delta_{g(t)}J+\langle X,\nabla J\rangle+(\beta-C_\tau)J.
\]
Since \(\beta-C_\tau\le 0\) on \(\mathcal M^n\times[0,\tau]\), the function \(J\) is a supersolution of the heat equation at points where \(J\le 0\). By Theorem~\ref{ch04-thm-scalar-max-principle-supersolution-lower-bound}, \(J\ge 0\) on \(\mathcal M^n\times[0,\tau]\), hence \(u\ge 0\) there as well. As \(\tau\) was arbitrary, the result follows.
\end{proof}

\subsection*{1.3. The heat equation with a nonlinear reaction term}

\begin{theorem}[Scalar maximum principle, fourth version]
\label{ch04-thm-scalar-max-principle-ode-bounds}
\source{chow-knopf-ricci-flow:page-0109}
\dcref{Theorem 4.4}
\group{weak-maximum-principles-scalar-equations}

Let \(u : \mathcal M^n \times [0,T) \to \mathbb R\) be a \(C^2\) supersolution to
\[
\frac{\partial v}{\partial t} = \Delta_{g(t)} v + \langle X, \nabla v \rangle + F(v)
\]
on a closed manifold. Suppose there exists \(C_1 \in \mathbb R\) such that \(u(x,0)\ge C_1\) for all \(x \in \mathcal M^n\), and let \(\varphi_1\) solve
\[
\frac{d\varphi_1}{dt}=F(\varphi_1),\qquad \varphi_1(0)=C_1.
\]
Then \(u(x,t)\ge \varphi_1(t)\) for all \(x \in \mathcal M^n\) and all \(t\) for which \(\varphi_1(t)\) exists.

Similarly, if \(u\) is a subsolution and \(u(x,0)\le C_2\) for all \(x \in \mathcal M^n\), and if \(\varphi_2\) solves
\[
\frac{d\varphi_2}{dt}=F(\varphi_2),\qquad \varphi_2(0)=C_2,
\]
then \(u(x,t)\le \varphi_2(t)\) for all \(x \in \mathcal M^n\) and all \(t\) for which \(\varphi_2(t)\) exists.
\end{theorem}

\begin{proof}
\sketch \uses{ch04-prop-scalar-max-principle-linear-reaction}
\source{chow-knopf-ricci-flow:page-0109}
We prove only the lower bound. The difference \(u-\varphi_1\) satisfies
\[
\frac{\partial}{\partial t}(u-\varphi_1)\ge \Delta(u-\varphi_1)+\langle X,\nabla(u-\varphi_1)\rangle+F(u)-F(\varphi_1).
\]
Fix \(\tau\in(0,T)\). Because \(\mathcal M^n\) is compact, there is \(C_\tau<\infty\) with \(|u|,|\varphi_1|\le C_\tau\) on \(\mathcal M^n\times[0,\tau]\). Since \(F\) is locally Lipschitz, there is \(L_\tau\) such that
\[
|F(v)-F(w)|\le L_\tau |v-w|
\]
on \([-C_\tau,C_\tau]\). Hence
\[
\frac{\partial}{\partial t}(u-\varphi_1)\ge \Delta(u-\varphi_1)+\langle X,\nabla(u-\varphi_1)\rangle -L_\tau \sgn(u-\varphi_1)\cdot(u-\varphi_1).
\]
Applying Proposition~\ref{ch04-prop-scalar-max-principle-linear-reaction} with \(\beta:=-L_\tau\sgn(u-\varphi_1)\) gives \(u-\varphi_1\ge 0\) on \(\mathcal M^n\times[0,\tau]\). Since \(\tau\) is arbitrary, the theorem follows.
\end{proof}

\begin{remark}
\label{ch04-rem-use-parabolic-maximum-principle}
\source{chow-knopf-ricci-flow:page-0109}
\dcref{Remark 4.5}
\group{weak-maximum-principles-scalar-equations}

In what follows, Theorem~\ref{ch04-thm-scalar-max-principle-ode-bounds} may be invoked simply as the parabolic maximum principle.
\end{remark}

\section*{2. Weak maximum principles for tensor equations}

\begin{theorem}[Tensor maximum principle, first version]
\label{ch04-thm-tensor-max-principle-nonnegativity}
\source{chow-knopf-ricci-flow:page-0110, chow-knopf-ricci-flow:page-0111, chow-knopf-ricci-flow:page-0112, chow-knopf-ricci-flow:page-0113}
\dcref{Theorem 4.6}
\group{weak-maximum-principles-tensor-equations}

Let \(g(t)\) be a smooth \(1\)-parameter family of Riemannian metrics on a closed manifold \(M^n\). Let \(\alpha(t)\in C^\infty(T^*M^n \symtensor T^*M^n)\) be a symmetric \((2,0)\)-tensor satisfying
\[
\frac{\partial}{\partial t}\alpha \ge \Delta_{g(t)}\alpha+\beta,
\]
where \(\beta(\alpha,g,t)\) is a symmetric \((2,0)\)-tensor that is locally Lipschitz in all variables and satisfies the null eigenvector assumption
\[
\beta(V,V)(x,t)\ge 0
\]
whenever \(V(x,t)\) is a null eigenvector of \(\alpha(t)\), i.e. whenever \((\alpha_{ij}V^j)(x,t)=0\). If \(\alpha(0)\ge 0\), then \(\alpha(t)\ge 0\) for all \(t\ge 0\) for which the solution exists.
\end{theorem}

\begin{proof}
\source{chow-knopf-ricci-flow:page-0110, chow-knopf-ricci-flow:page-0111, chow-knopf-ricci-flow:page-0112, chow-knopf-ricci-flow:page-0113}
\sketch The proof follows the scalar maximum principle by testing \(\alpha\) against a carefully chosen vector field. Suppose first that \(\alpha>0\) for \(0\le t<t_0\), and that at \((x_0,t_0)\) there is a vector \(v\) with \(\alpha_{ij}v^j(x_0,t_0)=0\). Extend \(v\) to a local space-time vector field \(V\) such that at \((x_0,t_0)\),
\[
\frac{\partial V}{\partial t}=0,\qquad \nabla V=0,\qquad \Delta V=0.
\]
Then the scalar function \(\alpha_{ij}V^iV^j\) satisfies
\[
\frac{\partial}{\partial t}(\alpha_{ij}V^iV^j)
 = \left(\frac{\partial}{\partial t}\alpha_{ij}\right)V^iV^j
 = (\Delta\alpha_{ij}+\beta_{ij})V^iV^j.
\]
At \((x_0,t_0)\), the minimum-point argument gives \(\Delta(\alpha_{ij}V^iV^j)\ge 0\), while the choice of \(V\) yields
\[
\Delta(\alpha_{ij}V^iV^j)=(\Delta\alpha_{ij})V^iV^j.
\]
Using the null eigenvector assumption then forces
\[
\frac{\partial}{\partial t}(\alpha_{ij}V^iV^j)\ge 0
\]
at \((x_0,t_0)\), contradicting the possibility that the first zero could decrease further.

For the full argument, fix \(\tau \in (0,T)\) and choose \(\delta>0\) so that for every \(t_0\in[0,\tau-\delta]\), positivity at time \(t_0\) propagates to \([t_0,t_0+\delta]\). For \(0<\varepsilon\le 1\), set
\[
A_\varepsilon(x,t):=\alpha(x,t)+\varepsilon[\delta+(t-t_0)]\,g(x,t).
\]
For \(\delta\) sufficiently small, \(A_\varepsilon\) is strictly positive definite at \(t=t_0\), and the evolution of \(A_\varepsilon\) satisfies
\[
\frac{\partial}{\partial t}A_\varepsilon
\ge \Delta A_\varepsilon+\beta(A_\varepsilon,g,t)+[\beta(\alpha,g,t)-\beta(A_\varepsilon,g,t)]+\varepsilon g+\varepsilon[\delta+(t-t_0)]\frac{\partial g}{\partial t}.
\]
Choosing \(\delta_0\) small enough that \(\partial_t g\ge -\frac1{4\delta_0}g\) on \([t_0,t_0+\delta_0]\) gives
\[
\varepsilon g+\varepsilon[\delta_0+(t-t_0)]\frac{\partial g}{\partial t}\ge \frac12 \varepsilon g.
\]
Local Lipschitz continuity of \(\beta\) gives a constant \(K\) with
\[
\beta(\alpha,g,t)-\beta(A_\varepsilon,g,t)\ge -K\varepsilon[\delta_0+(t-t_0)]g,
\]
and choosing \(\delta<1/(4K)\) yields
\[
\frac{\partial}{\partial t}A_\varepsilon>\Delta A_\varepsilon+\beta(A_\varepsilon,g,t).
\]
If \(A_\varepsilon\) were to lose positivity first at some \((x_1,t_1)\), choose a null eigenvector \(v\) there and extend it as above. Then
\[
0\ge \frac{\partial}{\partial t}\bigl((A_\varepsilon)_{ij}V^iV^j\bigr)
> \Delta\bigl((A_\varepsilon)_{ij}V^iV^j\bigr)+\beta_{ij}(A_\varepsilon,g,t)V^iV^j\ge 0,
\]
a contradiction. Thus \(A_\varepsilon>0\), and letting \(\varepsilon\searrow 0\) gives \(\alpha\ge 0\).
\end{proof}

\begin{remark}
\label{ch04-rem-tensor-maximum-principle-oversight}
\source{chow-knopf-ricci-flow:page-0113}
\dcref{Remark 4.7}
\group{weak-maximum-principles-tensor-equations}

The proof above corrects a minor oversight in the original argument of Section~9 of \([58]\), which omitted the \(\partial g/\partial t\) term.
\end{remark}

\section*{3. Advanced weak maximum principles for systems}

\begin{theorem}[Tensor maximum principle, second version]
\label{ch04-thm-tensor-max-principle-ode-pde}
\source{chow-knopf-ricci-flow:page-0114}
\dcref{Theorem 4.8}
\group{advanced-weak-maximum-principles-systems}

Let \(\mathcal M^n\) be a closed oriented manifold with a smooth family of metrics \(g(t)\) and Levi-Civita connections \(\nabla(t)\). Let \(\pi:\mathcal E\to\mathcal M^n\) be a vector bundle with bundle metric \(h\), and let \(\bar\nabla(t)\) be a smooth family of \(h\)-compatible connections. Let \(F:\mathcal E\times[0,T]\to\mathcal E\) be continuous, fiber-preserving, and fiberwise Lipschitz, and let \(\mathcal K\subset \mathcal E\) be closed. Suppose \(\alpha(t)\) solves
\[
\frac{\partial}{\partial t}\alpha=\hat\Delta\alpha+F(\alpha)
\]
with \(\alpha(0)\in\mathcal K\). Assume that \(\mathcal K\) is invariant under parallel translation by \(\bar\nabla(t)\) for all \(t\), each fiber slice \(\mathcal K_x\) is closed and convex, and every fiberwise ODE solution
\[
\frac{d}{dt}a=F(a),\qquad a(0)\in \mathcal K_x
\]
remains in \(\mathcal K_x\). Then \(\alpha(t)\in\mathcal K\) for all \(t\in[0,T]\).
\end{theorem}

\begin{theorem}[Tensor maximum principle, third version]
\label{ch04-thm-tensor-max-principle-time-dependent}
\source{chow-knopf-ricci-flow:page-0114, chow-knopf-ricci-flow:page-0115}
\dcref{Theorem 4.9}
\group{advanced-weak-maximum-principles-systems}

With the same setup, allow \(\mathcal K(t)\subset\mathcal E\) to vary with time. If \(\alpha(0)\in \mathcal K(0)\), the space-time track \(\bigcup_{t\in[0,T]}(\mathcal K(t)\times\{t\})\) is closed in \(\mathcal E\times[0,T]\), each \(\mathcal K(t)\) is invariant under parallel translation by \(\bar\nabla(t)\), each fiber slice \(\mathcal K_x(t)\) is closed and convex, and every fiberwise ODE solution starting in \(\mathcal K_x(0)\) remains in \(\mathcal K_x(t)\), then \(\alpha(t)\in\mathcal K(t)\) for all \(t\in[0,T]\).
\end{theorem}

\begin{theorem}[Tensor maximum principle, fourth version]
\label{ch04-thm-tensor-max-principle-avoidance}
\source{chow-knopf-ricci-flow:page-0115}
\dcref{Theorem 4.10}
\group{advanced-weak-maximum-principles-systems}

With the same setup, let \(\mathcal A(t)\subset\mathcal K(t)\) be an avoidance set such that the space-time tracks of \(\mathcal K(t)\) and \(\mathcal A(t)\) are closed in \(\mathcal E\times[0,T]\), each \(\mathcal K(t)\) is invariant under parallel translation by \(\bar\nabla(t)\), each \(\mathcal K_x(t)\) is closed and convex, and \(\alpha(t)\) avoids \(\mathcal A(t)\) for all \(t\in[0,T]\). If every fiberwise ODE solution starting in \(\mathcal K_x(0)\setminus\mathcal A_x(0)\) either remains in \(\mathcal K_x(t)\) for all \(t\in[0,T]\) or enters \(\mathcal A(t_0)\) at some \(t_0\in(0,T]\), then \(\alpha(t)\in\mathcal K(t)\) for all \(t\in[0,T]\).
\end{theorem}
